% =====================================================================
%  Άσκηση 4 - Συγχρονισμός, εκτίμηση CFO και καναλιού
%  Ασύρματες Επικοινωνίες, Πολυτεχνείο Κρήτης
%  Μεταγλώττιση:  xelatex Wireless_Communications_EX4_report.tex  (δύο φορές)
% =====================================================================
\documentclass[11pt,a4paper]{article}

\usepackage{fontspec}
\usepackage{polyglossia}
\setmainlanguage{greek}
\setotherlanguage{english}
\setmainfont{Libertinus Serif}
\setsansfont{Libertinus Sans}
\setmonofont{Libertinus Mono}
\newfontfamily\greekfont{Libertinus Serif}
\newfontfamily\greekfontsf{Libertinus Sans}
\newfontfamily\greekfonttt{Libertinus Mono}
\newfontfamily\englishfont{Libertinus Serif}
\newfontfamily\englishfontsf{Libertinus Sans}
\newfontfamily\englishfonttt{Libertinus Mono}

\usepackage{amsmath,amssymb}
\usepackage{graphicx}
\usepackage[margin=2.5cm]{geometry}
\usepackage{float}
\usepackage{caption}
\usepackage{booktabs}

\graphicspath{{figures/}}
\captionsetup{font=normalsize}
\setlength{\parskip}{4pt}
\emergencystretch=2em

% Το ζητούμενο της εκφώνησης με έντονα γράμματα (όπως στις προηγούμενες αναφορές)
\newcommand{\zit}[1]{\par\noindent\textbf{#1}\par\medskip}

\newcommand{\ovr}{\mathrm{over}}
\newcommand{\Ntr}{N_{\mathrm{tr}}}

% Δύο σχήματα δίπλα-δίπλα
\newcommand{\twofig}[4]{%
\begin{figure}[H]
  \centering
  \begin{minipage}[t]{0.49\textwidth}\centering
    \includegraphics[width=\linewidth]{#1}
    \caption{#2}
  \end{minipage}\hfill
  \begin{minipage}[t]{0.49\textwidth}\centering
    \includegraphics[width=\linewidth]{#3}
    \caption{#4}
  \end{minipage}
\end{figure}}

\newcommand{\onefig}[3][0.72]{%
\begin{figure}[H]
  \centering
  \includegraphics[width=#1\textwidth]{#2}
  \caption{#3}
\end{figure}}

\begin{document}

% ---------------------------------------------------------------------
\begin{center}
  {\large \textenglish{Technical University of Crete}}\\[2pt]
  \textenglish{School of Electrical and Computer Engineering}\\[14pt]
  {\LARGE\bfseries Άσκηση 4 - Συγχρονισμός, CFO και Εκτίμηση Καναλιού}\\[8pt]
  Ασύρματες Επικοινωνίες\\[14pt]
  \begin{tabular}{ll}
    \textenglish{Student:}    & \textenglish{Emmanouil Thomas Chatzakis} \\
    \textenglish{Student ID:} & 2021030061
  \end{tabular}
\end{center}
\vspace{6pt}
\hrule
\vspace{14pt}

Σε αυτή την άσκηση μελετάμε τεχνικές συγχρονισμού, εκτίμησης CFO και εκτίμησης καναλιού σε
ιδανικά και μη-ιδανικά κανάλια. Η θεωρία που χρησιμοποιούμε είναι από τις σημειώσεις του
μαθήματος: \emph{Κεφάλαιο 4 – Εκτίμηση και Ισοστάθμιση} (εκτίμηση ελαχίστων τετραγώνων, ZF
ισοσταθμιστής), \emph{Κεφάλαιο 5 – Συγχρονισμός, Διακριτό ισοδύναμο κανάλι} και \emph{Κεφάλαιο 6 –
Εκτίμηση και διόρθωση CFO}. Οι αριθμοί εξισώσεων που αναφέρουμε είναι αυτοί των σημειώσεων.

Η υλοποίηση έγινε σε MATLAB. Το κύριο αρχείο είναι το \texttt{Ex4.m} και κάθε βασικό βήμα της
θεωρίας είναι σε ξεχωριστή συνάρτηση:
\begin{itemize}
  \item \texttt{energy\_sync.m}: ενέργεια των symbol-spaced υπακολουθιών $\mathrm{E}_d$,
  \item \texttt{training\_corr.m}: συσχέτιση με τα σύμβολα εκπαίδευσης $\mathcal{C}_d$, (5.18),
  \item \texttt{cfo\_estimate\_fft.m}: εκτίμηση CFO με μετασχηματισμό Fourier, Κεφ. 6,
  \item \texttt{training\_matrix.m}, \texttt{ls\_channel\_estimate.m}: εκτίμηση καναλιού LS, (4.9),
  \item \texttt{zf\_equalizer.m}: ZF ισοσταθμιστής, (4.17)–(4.18),
  \item \texttt{joint\_cfo\_channel\_ls.m}: ταυτόχρονη εκτίμηση CFO και καναλιού, (5.34)–(5.41).
\end{itemize}

Σε όλη την άσκηση χρησιμοποιούμε $T=1$, $\ovr = T/T_s = 10$, SRRC φίλτρα με $\beta = 0.4$ και
μισό μήκος $A = 3$ περιόδους συμβόλου, $N = 200$ σύμβολα ανά πακέτο και $\Ntr = 50$ σύμβολα
εκπαίδευσης. Τα κανάλια είναι αθόρυβα. Για να είναι τα αποτελέσματα αναπαραγώγιμα
χρησιμοποιούμε \texttt{rng(2021030061)} και όλες οι αριθμητικές τιμές της αναφοράς βρίσκονται
στο αρχείο \texttt{results\_log.txt}.

% =====================================================================
\section{Πρώτο Μέρος}

\zit{Αρχικά, θα ασχοληθούμε με ιδανικά κανάλια.}

\subsection*{1.1}
\zit{Να δημιουργήσετε πακέτο με 4-QAM σύμβολα, $A_0,\dots,A_{N-1}$. Να θεωρήσετε τα πρώτα $\Ntr$
σύμβολα ως σύμβολα εκπαίδευσης. Να περάσετε το πακέτο από το φίλτρο μορφοποίησης και να
υπολογίσετε το χαμηλοπερατό ισοδύναμο της εισόδου του καναλιού.}

Δημιουργούμε $N=200$ ισοπίθανα σύμβολα 4-QAM με ενέργεια συμβόλου
\[
A_n \in \{\pm 1 \pm j\}, \qquad \mathbb{E}\left[|A_n|^2\right] = \sigma_A^2 = 2 ,
\]
και θεωρούμε τα πρώτα $\Ntr = 50$ ως σύμβολα εκπαίδευσης. Τοποθετούμε το σύμβολο $A_n$ κάθε
$\ovr$ δείγματα (στη χρονική στιγμή $nT$) και κάνουμε συνέλιξη με τον δειγματοληπτημένο παλμό
SRRC, ο οποίος έχει μετατοπιστεί ώστε να είναι αιτιατός. Έτσι λαμβάνουμε τα δείγματα
\[
A(kT_s) = \sum_{n=0}^{N-1} A_n\, g_T(kT_s - nT).
\]

\onefig[0.7]{p1_At}{Low-pass equivalent channel input A(t)}

\subsection*{1.2}
\zit{Να θεωρήσετε ότι η κρουστική απόκριση του φυσικού καναλιού είναι $c(t) = \delta(t)$ και ότι το
κανάλι είναι αθόρυβο.}

Η έξοδος του καναλιού είναι απλά $A(t) * \delta(t) = A(t)$.

\subsection*{1.3}
\zit{Να φιλτράρετε με το προσαρμοσμένο φίλτρο στο δέκτη.}

Η έξοδος του προσαρμοσμένου φίλτρου υπολογίζεται ως
\[
Y_k = Y(kT_s) = T_s \sum_m A(mT_s)\, g_R(kT_s - mT_s),
\]
όπου ο παράγοντας $T_s$ προσεγγίζει το ολοκλήρωμα της συνέλιξης συνεχούς χρόνου. Ο χρόνος στην
έξοδο αρχίζει στο $t = 0$.

\subsection*{1.4}
\zit{Να υπολογίσετε το σύνθετο αναλογικό κανάλι, υποθέτοντας ότι ο χρόνος αρχίζει στο $t=0$. Να
σχεδιάσετε την απόλυτη τιμή του.}

Το σύνθετο αναλογικό κανάλι (5.5) είναι
\[
h(t) = g_T(t) * c(t) * g_R(t) = g_T(t) * g_R(t),
\]
δηλαδή ένας (περικομμένος) παλμός raised cosine. Επειδή η ενέργεια του SRRC είναι μονάδα, η τιμή
στο μέγιστο είναι $|h(6T)| = 0.9996 \approx 1$.

\onefig[0.7]{p1_composite_c1}{Composite analog channel |h(t)|, c = 1}

Παρατηρούμε ότι το κανάλι έχει διάρκεια $4AT = 12T$ και μέγιστο στο $t = 2AT = 6T$ (δείγμα 60),
λόγω της μετατόπισης που το κάνει αιτιατό. Επίσης μηδενίζεται στα σημεία $6T \pm kT$, αφού είναι
παλμός Nyquist.

\subsection*{1.5}
\zit{Να προσπαθήσετε να συγχρονιστείτε με βάση την ενέργεια του ληφθέντος σήματος, θεωρώντας ότι
$d = 0,\dots,4A\ovr-1$. Να σχεδιάσετε το κατάλληλο στατιστικό. Τι παρατηρείτε;}

Σύμφωνα με το εδάφιο 5.1.1, για κάθε καθυστέρηση $d$ υπολογίζουμε την ενέργεια της symbol-spaced
υπακολουθίας μήκους $N$
\[
\mathrm{E}_d = \sum_{n=0}^{N-1} |Y_{d + n\,\ovr}|^2, \qquad
\boxed{d^* = \arg\max_{d} \mathrm{E}_d}.
\]

\onefig[0.7]{p1_Ed_c1}{Energy statistic $\mathrm{E}_d$}

Βρίσκουμε $d^* = 60$, δηλαδή ακριβώς τη θέση του μεγίστου του σύνθετου καναλιού. Η μέγιστη τιμή
είναι $\mathrm{E}_{60} = 399.27 \approx N\sigma_A^2 = 400$.

Παρατηρούμε όμως ότι το στατιστικό είναι περιοδικό με περίοδο $\ovr = 10$ δείγματα. Τα τοπικά
μέγιστα στα $d = 60 \pm 10k$ αντιστοιχούν στη σωστή φάση δειγματοληψίας αλλά σε λάθος αρχή του
πακέτου κατά $k$ σύμβολα. Η δεύτερη μεγαλύτερη τιμή είναι $397.26$, δηλαδή μικρότερη μόνο κατά
$\approx |A_n|^2 = 2$, αφού χάνεται ένα μόνο σύμβολο. Άρα η μέθοδος της ενέργειας βρίσκει
εύκολα τη φάση δειγματοληψίας, αλλά η αρχή του πακέτου δεν ξεχωρίζει καθαρά και σε ενθόρυβο
περιβάλλον θα μπορούσε εύκολα να βρεθεί λάθος.

\subsection*{1.6}
\zit{Να προσπαθήσετε να συγχρονιστείτε με βάση τα σύμβολα εκπαίδευσης, για τα ίδια $d$. Να
σχεδιάσετε την απόλυτη τιμή του αντίστοιχου στατιστικού και να τη συγκρίνετε με την απόλυτη τιμή
του σύνθετου καναλιού. Τι παρατηρείτε;}

Σύμφωνα με το εδάφιο 5.2.2 υπολογίζουμε
\[
\mathcal{C}_d = \sum_{n=0}^{\Ntr-1} A_n^*\, Y_{d + n\,\ovr}.
\]
Αφού τα $A_n$ είναι ανεξάρτητα με μηδενική μέση τιμή, από την (5.19)
$\mathbb{E}[\mathcal{C}_d] = \Ntr \sigma_A^2\, h(dT_s)$ και άρα
\[
\boxed{\hat h(dT_s) = \frac{\mathcal{C}_d}{\Ntr\sigma_A^2}}.
\]

\onefig[0.7]{p1_Cd_c1}{Training statistic vs composite channel}

Το μέγιστο του $|\mathcal{C}_d|$ είναι πάλι στο $d^* = 60$. Παρατηρούμε ότι γύρω από τον κύριο
λοβό το $|\hat h_d|$ ταυτίζεται με το $|h(dT_s)|$ και έχει ένα μοναδικό, καθαρό μέγιστο. Σε
αντίθεση με την ενέργεια, εδώ δεν υπάρχει αμφιβολία για την αρχή του πακέτου, αφού τα σύμβολα
εκπαίδευσης είναι γνωστά στον δέκτη.

Μακριά από τον κύριο λοβό το $|\hat h_d|$ είναι περίπου $0.2$–$0.3$ αντί για μηδέν. Αυτό
οφείλεται στους όρους $A_n^* A_k h^d_{n-k}$ για $n \neq k$ της (5.18). Έχουν μηδενική μέση τιμή,
αλλά δεν μηδενίζονται για πεπερασμένο $\Ntr = 50$, οπότε $\mathcal{C}_d \approx \mathbb{E}[\mathcal{C}_d]$
μόνο κατά προσέγγιση.

\subsection*{1.7}
\zit{Αφού συγχρονιστείτε με έναν από τους δύο παραπάνω τρόπους, να κατασκευάσετε την μήκους $N$
symbol-spaced ακολουθία εξόδου και να τη σχεδιάσετε. Τι παρατηρείτε;}

Η ακολουθία εξόδου είναι $r_k = Y_{d^* + k\,\ovr}$, $k = 0,\dots,N-1$.

\onefig[0.62]{p1_output_c1}{Symbol-spaced output, c = 1}

Τα δείγματα πέφτουν πάνω στα σημεία του αστερισμού 4-QAM, αφού το $h(t)$ είναι παλμός Nyquist και
δειγματοληπτούμε στις σωστές στιγμές (δεν υπάρχει ISI). Η μέγιστη απόκλιση είναι
$\max_k |r_k - A_k| = 3.2\cdot 10^{-2}$ και οφείλεται στην περικοπή του SRRC σε $\pm 3T$, λόγω της
οποίας το $g_T * g_R$ δεν είναι ακριβώς Nyquist.

\subsection*{1.8}
\zit{Να επαναλάβετε τα παραπάνω με τη μόνη διαφορά στο ότι $c(t) = c\,\delta(t)$, με $c \in \mathbb{C}$.
Τι παρατηρείτε;}

Επιλέγουμε τυχαία $c \sim \mathcal{CN}(0,1)$ και προκύπτει $c = -0.7093 + 0.6824j$, με
$|c| = 0.9843$ και $\angle c = 136.11^\circ$. Τώρα $h(t) = c\, g_T(t) * g_R(t)$.

\twofig{p1_Ed_crand}{Energy statistic, complex c}{p1_Cd_crand}{Training statistic, complex c}

Και οι δύο μέθοδοι δίνουν πάλι $d^* = 60$. Το $\mathrm{E}_d$ απλά πολλαπλασιάζεται με $|c|^2$
($386.82 = 0.969 \cdot 399.27$) και το $|\hat h_d|$ με $|c|$, άρα η θέση του μεγίστου δεν
αλλάζει.

\onefig[0.62]{p1_output_crand}{Symbol-spaced output, complex c}

Η ακολουθία εξόδου είναι $r_k = c\,A_k$. Ο αστερισμός είναι στραμμένος κατά $\angle c = 136^\circ$
και κλιμακωμένος κατά $|c|$, οπότε τα σημεία βρίσκονται περίπου πάνω στους άξονες. Ο συγχρονισμός
δεν επηρεάζεται, αλλά για να αποφασίσουμε σωστά τα σύμβολα πρέπει πρώτα να εκτιμήσουμε το $c$
(coherent detection).

% =====================================================================
\section{Δεύτερο Μέρος}

\zit{Στο δεύτερο μέρος, στο ιδανικό κανάλι θα προσθέσουμε CFO.}

\subsection*{2.1 – 2.2}
\zit{Να επαναλάβετε τα βήματα 1.1 και 1.2. Να υποθέσετε την ύπαρξη CFO $\Delta F$ και να
δειγματοληπτήσετε, λαμβάνοντας $Y_n = e^{j(2\pi(\Delta F T_s)n + \phi)} A(nT_s)$, με
$\Delta F_s := \Delta F T_s \approx 10^{-3}$ και $\phi = 0$.}

Χρησιμοποιούμε το ίδιο πακέτο και $c = 1$. Το σήμα στην είσοδο του δέκτη είναι
\[
Y_n = e^{j(2\pi \Delta F_s n + \phi)} A(nT_s), \qquad \Delta F_s = 10^{-3},\; \phi = 0 .
\]
Το CFO ανά σύμβολο είναι $\Delta f = \Delta F T = \ovr \cdot \Delta F_s = 0.01$, δηλαδή η φάση
αλλάζει κατά $3.6^\circ$ ανά σύμβολο. Σε όλο το πακέτο των $N = 200$ συμβόλων η φάση κάνει
$2\pi \Delta f N = 4\pi$, δηλαδή δύο πλήρεις περιστροφές.

\subsection*{2.3}
\zit{Να επαναλάβετε τα βήματα 1.3 μέχρι 1.7. Τι παρατηρείτε;}

Αφού $\Delta F \ll 1/T$, από τις (6.10)–(6.11) και (5.25) η έξοδος του προσαρμοσμένου φίλτρου
προσεγγίζεται από
\[
Y(t) \approx e^{j(2\pi\Delta F t + \phi)} \sum_{n=0}^{N-1} A_n\, h(t - nT).
\]
Το σύνθετο κανάλι είναι το ίδιο με το πρώτο μέρος.

\twofig{p2_Ed_a}{Energy statistic with CFO}{p2_Cd_a}{Training statistic with CFO}

\textbf{Ενέργεια.} Κάθε δείγμα πολλαπλασιάζεται με εκθετικό μέτρου 1, άρα το $\mathrm{E}_d$ δεν
αλλάζει και βρίσκουμε πάλι $d^* = 60$.

\textbf{Σύμβολα εκπαίδευσης.} Από τις (5.26)–(5.28),
$\mathbb{E}[\mathcal{C}_d] = \alpha_d\, h(dT_s)$ με
\[
|\alpha_d| = \sigma_A^2 \left| \sum_{n=0}^{\Ntr-1} e^{j2\pi\Delta f n} \right|
= \sigma_A^2 \left| \frac{\sin(\pi \Delta f \Ntr)}{\sin(\pi\Delta f)} \right| .
\]
Για $\Delta f = 0.01$ και $\Ntr = 50$ η θεωρία δίνει $|\alpha|/(\Ntr\sigma_A^2) = 0.6367$ και
στην προσομοίωση μετράμε $\max|\hat h_d| / \max|h| = 0.6367$, δηλαδή ακριβώς την ίδια τιμή.
Το στατιστικό εξασθενεί, αλλά η μορφή του κύριου λοβού παραμένει. Αν όμως το $\Delta f\,\Ntr$ ήταν
ακέραιος (π.χ. $\Ntr = 100$), τότε $\alpha = 0$ και ο συγχρονισμός με σύμβολα εκπαίδευσης θα
αποτύγχανε εντελώς.

Εδώ το μέγιστο του $|\mathcal{C}_d|$ βρέθηκε στο $d = 61$ αντί για $60$. Ο κύριος λοβός είναι
πολύ επίπεδος κοντά στο μέγιστο και, μετά την εξασθένηση κατά $|\alpha|$, οι όροι $n \neq k$ τον
μετακινούν κατά ένα δείγμα. Γι' αυτό στο δεύτερο μέρος κατασκευάζουμε την ακολουθία εξόδου με
τον συγχρονισμό ενέργειας ($d^* = 60$), ο οποίος δεν επηρεάζεται από το CFO.

\onefig[0.62]{p2_output_a}{Symbol-spaced output with CFO (no correction)}

Η ακολουθία εξόδου είναι $r_k \approx A_k\, e^{j(2\pi\Delta f k + \phi'')}$. Επειδή η φάση αλλάζει
συνεχώς, τα δείγματα απλώνονται σε κύκλο ακτίνας $\sqrt 2$ και η απόφαση για τα σύμβολα είναι
αδύνατη χωρίς διόρθωση του CFO.

\subsection*{2.4}
\zit{Να υλοποιήσετε και το βήμα 1.8 ή/και να θέσετε τυχαία φάση $\Phi \sim \mathcal{U}[0, 2\pi)$.
Τι παρατηρείτε;}

Υλοποιούμε και τα δύο μαζί: το ίδιο $c = -0.7093 + 0.6824j$ με το 1.8 και τυχαία φάση
$\Phi = 5.4516$ rad.

\onefig[0.62]{p2_output_b}{Symbol-spaced output with CFO, complex c and random phase}

Οι δύο συγχρονισμοί δίνουν ακριβώς τα ίδια $d$ με την περίπτωση $c=1,\ \phi=0$, αφού το $c$ και το
$\Phi$ αλλάζουν μόνο το κοινό πλάτος και την κοινή φάση όλων των δειγμάτων. Η ακολουθία εξόδου
είναι πάλι κύκλος, με ακτίνα $\sqrt 2 |c|$.

\subsection*{2.5}
\zit{Αν έχετε καταφέρει να συγχρονιστείτε, να εκτιμήσετε το CFO με χρήση συμβόλων εκπαίδευσης.}

Ακολουθούμε τη διαδικασία του Κεφαλαίου 6. Μετά τον συγχρονισμό
$r_k \approx \tilde h\, A_k\, e^{j2\pi\Delta f k}$, όπου το $\tilde h$ περιέχει το $c$ και όλες τις
σταθερές φάσεις. Πολλαπλασιάζουμε τα δείγματα που αντιστοιχούν στα σύμβολα εκπαίδευσης με
$A_l^*$ (6.16):
\[
z_l = r_l A_l^* = |A_l|^2 \tilde h\, e^{j2\pi \Delta f\, l}, \qquad l = 0,\dots,\Ntr-1 .
\]
Η $\{z_l\}$ είναι μιγαδικό εκθετικό συχνότητας $\Delta f$ (με το πρόσημο της (5.25); στην (6.16)
το πρόσημο είναι αρνητικό λόγω του τρόπου αποδιαμόρφωσης). Η εκτίμηση είναι το σημείο στο οποίο
μεγιστοποιείται το μέτρο του μετασχηματισμού Fourier,
\[
\boxed{\Delta\hat f = \arg\max_f \left| \mathcal{F}\{z_l\}(f) \right|},
\]
τον οποίο υπολογίζουμε με FFT $2^{16}$ σημείων (zero padding) για λεπτό πλέγμα συχνοτήτων.

\twofig{p2_cfo_fft_a}{Fourier transform of $z_l$ (case c = 1)}{p2_cfo_fft_b}{Fourier transform of $z_l$ (complex c, random phase)}

Και στις δύο περιπτώσεις $\Delta\hat f = 0.010010$ (πραγματικό $0.01$), δηλαδή
$\Delta \hat F_s = 1.001 \cdot 10^{-3}$. Το $c$ και η φάση αλλάζουν μόνο τη φάση της $z_l$ και όχι
τη συχνότητά της. Το $|\mathcal{F}\{z_l\}|$ έχει τη μορφή του πυρήνα Dirichlet, με μηδενισμούς
ανά $1/\Ntr = 0.02$.

Αν είχαμε κατασκευάσει την ακολουθία από το $d = 61$ του $|\mathcal{C}_d|$, η εκτίμηση θα ήταν
$\Delta\hat f = 0.009857$ και η τελική μέγιστη απόκλιση από τα σύμβολα $0.519$ (αντί για $0.036$).
Το σφάλμα χρονισμού $0.1T$ δημιουργεί ISI, η οποία πολώνει την εκτίμηση του CFO.

\subsection*{2.6}
\zit{Να αναιρέσετε την επίδραση του CFO και να σχεδιάσετε την μήκους $N$ symbol-spaced ακολουθία
εξόδου. Τι παρατηρείτε;}

Από την (6.17) η διορθωμένη ακολουθία είναι
\[
\tilde r_k = r_k\, e^{-j2\pi \Delta\hat f k} = \tilde h\, A_k, \qquad k = 0,\dots,N-1 .
\]

\twofig{p2_cfo_corrected_a}{After CFO correction (c = 1)}{p2_cfo_corrected_b}{After CFO correction (complex c, random phase)}

Ο κύκλος «μαζεύεται» ξανά σε τέσσερα σημεία, τα οποία όμως είναι στραμμένα κατά μία κοινή φάση:
περίπου $10.7^\circ$ για $c=1$ και $99.2^\circ$ στη δεύτερη περίπτωση. Η φάση αυτή είναι η $\phi''$
του Κεφαλαίου 6. Για $c=1,\ \phi = 0$ ισούται με τη φάση του CFO μέχρι το μέγιστο του πρώτου
παλμού (δείγμα $A\,\ovr = 30$ της $A(t)$): $2\pi \Delta F_s \cdot 30 = 10.8^\circ$. Στη δεύτερη
περίπτωση προστίθενται και τα $\angle c + \Phi$: $136.1^\circ + 312.4^\circ + 10.8^\circ = 459.3^\circ
\equiv 99.3^\circ$.

\subsection*{2.7}
\zit{Να εκτιμήσετε το διακριτό ισοδύναμο κανάλι (στην περίπτωση αυτή το κανάλι έχει μόνο ένα
συντελεστή) με τη μέθοδο των ελαχίστων τετραγώνων.}

Για τα σύμβολα εκπαίδευσης $\tilde{\mathbf r} = \mathbf a\, h + \mathbf v$, με
$\mathbf a = [A_0 \cdots A_{\Ntr-1}]^T$. Η λύση ελαχίστων τετραγώνων είναι
\[
\boxed{\hat h_{\mathrm{LS}} = (\mathbf a^H \mathbf a)^{-1} \mathbf a^H \tilde{\mathbf r}
 = \frac{\sum_l A_l^* \tilde r_l}{\sum_l |A_l|^2}},
\]
η οποία, αφού $|A_l|^2 = 2$ σταθερό, ταυτίζεται με την εκτίμηση (6.19) του Κεφαλαίου 6.

\begin{table}[H]
\centering
\begin{tabular}{lccc}
\toprule
Περίπτωση & $\hat h_{\mathrm{LS}}$ & $|\hat h_{\mathrm{LS}}|$ & $\angle \hat h_{\mathrm{LS}}$ \\
\midrule
$c = 1,\ \phi = 0$          & $0.9811 + 0.1860j$  & $0.9985$ & $10.73^\circ$ \\
$c$ μιγαδικό, $\Phi$ τυχαία & $-0.1571 + 0.9702j$ & $0.9828$ & $99.20^\circ$ \\
\bottomrule
\end{tabular}
\end{table}

Το μέτρο είναι $|c|\,h(6T)$ ($0.9843\cdot 0.9996 = 0.984$ στη δεύτερη περίπτωση) και η φάση είναι η
κοινή φάση που είδαμε στο 2.6.

\subsection*{2.8}
\zit{Να προσπαθήσετε να αναιρέσετε την επίδραση του καναλιού και να σχεδιάσετε την ακολουθία που
λαμβάνετε μετά από αυτό. Τι παρατηρείτε;}

Διαιρούμε με την εκτίμηση του καναλιού: $\hat A_k = \tilde r_k / \hat h_{\mathrm{LS}}$.

\twofig{p2_equalized_a}{Output after CFO and channel correction (c = 1)}{p2_equalized_b}{Output after CFO and channel correction (complex c, random phase)}

Τα δείγματα πέφτουν πάνω στα σημεία του αστερισμού και δεν έχουμε κανένα λάθος συμβόλου
(0/200) και στις δύο περιπτώσεις, με μέγιστη απόκλιση $0.036$. Η μικρή απόκλιση οφείλεται στην
περικοπή του SRRC και στην προσέγγιση (6.11), αφού η φάση του CFO αλλάζει λίγο κατά τη διάρκεια
κάθε παλμού. Το σφάλμα εκτίμησης $10^{-5}$ του $\Delta f$ δίνει σφάλμα φάσης μόλις $0.7^\circ$ σε
200 σύμβολα.

% =====================================================================
\section{Τρίτο Μέρος}

\zit{Στο τρίτο μέρος, θα ασχοληθούμε με μη-ιδανικά κανάλια χωρίς CFO.}

\subsection*{3.1 – 3.2}
\zit{Να υλοποιήσετε το βήμα 1.1. Να θεωρήσετε ότι η κρουστική απόκριση του φυσικού καναλιού
είναι $c(t) = c_0\delta(t) + c_1\delta(t - KT_s)$, με $K \sim \mathcal{U}\{1,\dots,4\ovr\}$.}

Χρησιμοποιούμε το ίδιο πακέτο με το πρώτο μέρος. Από την κλήρωση προέκυψε $K = 32$ δείγματα
($3.2T$) και, με $c_0, c_1 \sim \mathcal{CN}(0,1)$,
\[
c_0 = -0.1904 - 1.0210j, \qquad c_1 = 2.1523 + 0.0471j .
\]
Σε διακριτό χρόνο το φυσικό κανάλι είναι το διάνυσμα $[c_0,\, 0, \dots, 0,\, c_1]$ μήκους $K+1$.

\subsection*{3.3}
\zit{Να υλοποιήσετε τα βήματα 1.3 και 1.4.}

Το σύνθετο αναλογικό κανάλι είναι
\[
h(t) = g_T(t) * c(t) * g_R(t) = c_0\, h_{\mathrm{RC}}(t) + c_1\, h_{\mathrm{RC}}(t - KT_s),
\]
όπου $h_{\mathrm{RC}} = g_T * g_R$, με συνολική διάρκεια $4AT + KT_s = 15.2T$.

\onefig[0.7]{p3_composite}{Composite analog channel |h(t)|, K = 32}

Παρατηρούμε δύο κύριους λοβούς, στο $6T$ και στο $9.2T$. Ο δεύτερος είναι μεγαλύτερος αφού
$|c_1| > |c_0|$. Επειδή η απόσταση των λοβών δεν είναι ακέραιο πολλαπλάσιο του $T$, δεν υπάρχει
χρονική στιγμή δειγματοληψίας χωρίς ISI.

\subsection*{3.4}
\zit{Να υλοποιήσετε το βήμα 1.6, υπολογίζοντας την ποσότητα $\mathrm{corr}_d$, για
$d = 0,\dots,(4A\ovr + 1) + (K + 2) - 2$.}

Υπολογίζουμε το $\mathcal{C}_d$ για $d = 0,\dots,153$ και την εκτίμηση
$\hat h_d = \mathcal{C}_d/(\Ntr\sigma_A^2)$.

\onefig[0.7]{p3_Cd}{Training statistic vs composite channel (non-ideal channel)}

Το $|\hat h_d|$ ακολουθεί τους δύο λοβούς του $|h(dT_s)|$ και έχει μέγιστο στο $d = 92$
($9.2T$). Οι αποκλίσεις όμως είναι μεγαλύτερες από το πρώτο μέρος (π.χ. $1.38$ αντί για $1.04$ στον
πρώτο λοβό). Τώρα οι όροι $A_n^*A_k h^d_{n-k}$, $n\ne k$, περιέχουν τις ISI τιμές ενός καναλιού με
πολύ περισσότερη ενέργεια, οπότε με $\Ntr = 50$ η προσέγγιση
$\mathcal{C}_d \approx \mathbb{E}[\mathcal{C}_d]$ είναι χειρότερη.

\subsection*{3.5}
\zit{Να υπολογίσετε την ποσότητα $\mathrm{E}_d$, υποθέτοντας μήκος ισοδύναμου διακριτού καναλιού
$M$ (να βρείτε το κατάλληλο διάστημα για τις τιμές του $d$), και να συγχρονιστείτε, λαμβάνοντας μία
symbol-spaced ακολουθία μήκους $N + M - 1$.}

Επιλέγουμε $M = 10$. Το σημαντικό κομμάτι του $h_{\mathrm{RC}}$ είναι περίπου $4T$ ($\pm 2T$) και
η καθυστέρηση $KT_s$ προσθέτει έως $4T$, οπότε το σημαντικό μέρος του καναλιού είναι περίπου
$8T$. Σύμφωνα με τις σημειώσεις, μια μικρή υπερεκτίμηση του $M$ δεν έχει δραματικές συνέπειες,
ενώ η υποεκτίμηση μπορεί να υποβαθμίσει σημαντικά την απόδοση.

Από την (5.21)
\[
\mathrm{E}_d = \sum_{m=0}^{M-1} |\hat h_{d + m\,\ovr}|^2 .
\]
Για να υπάρχουν όλα τα δείγματα του $\hat h$ πρέπει $d + (M-1)\ovr \le d_{\max} = 153$, άρα
$d = 0,\dots,63$. Στο ίδιο σχήμα σχεδιάζουμε και το ίδιο στατιστικό με το πραγματικό $h$.

\onefig[0.7]{p3_Ed}{Channel energy $\mathrm{E}_d$ for M = 10}

Το μέγιστο είναι στο $d^* = 32$, όπως και με το πραγματικό κανάλι. Παρατηρούμε ότι το
$\mathrm{E}_d$ του πραγματικού καναλιού έχει ουσιαστικά ίσα μέγιστα ($5.6565$) σε όλα τα
$d = 2, 12, 22, 32, 42$. Αφού $M = 10$ σύμβολα καλύπτουν όλο το σημαντικό μέρος του καναλιού, κάθε
παράθυρο με τη σωστή φάση δειγματοληψίας ($d \bmod 10 = 2$) «μαζεύει» όλη την ενέργεια. Άρα η
φάση δειγματοληψίας βρίσκεται μονοσήμαντα, ενώ η θέση του παραθύρου επιλέγεται από τα μικρά σφάλματα
του $\hat h$. Η φάση $2$ αντιστοιχεί στη δειγματοληψία του ισχυρού λοβού στο $9.2T$.

Η ακολουθία εξόδου είναι (5.24) $Y_k = Y_{d^* + k\,\ovr}$, $k = 0,\dots,N+M-2$, και το
πραγματικό διακριτό ισοδύναμο κανάλι είναι (5.11) $h_m = h(d^*T_s + mT)$, $m = 0,\dots,M-1$.

\twofig{p3_discrete_channel}{Composite analog and discrete equivalent channel}{p3_output}{Symbol-spaced output of the non-ideal channel}

Στην ακολουθία εξόδου φαίνεται καθαρά η ISI. Οι δύο ισχυροί συντελεστές $h_3 \approx -0.19-0.95j$
και $h_6 \approx 2.15$ δίνουν $Y_k \approx h_6A_{k-6} + h_3A_{k-3}$, δηλαδή 16 ομάδες σημείων, οι
οποίες απλώνονται λίγο από τους μικρότερους συντελεστές. Τα λίγα σημεία κοντά στο μηδέν είναι τα
τελευταία δείγματα, όπου η έξοδος «σβήνει».

\subsection*{3.5 (β)}
\zit{Να εκτιμήσετε το διακριτό ισοδύναμο κανάλι μήκους $M$, με τη μέθοδο των ελαχίστων
τετραγώνων.}

Όπως στην άσκηση 3, οι έξοδοι που εξαρτώνται μόνο από σύμβολα εκπαίδευσης είναι οι
$Y_k$, $k = M-1,\dots,\Ntr-1$ (41 εξισώσεις):
\[
\begin{bmatrix} Y_{M-1} \\ \vdots \\ Y_{\Ntr-1} \end{bmatrix} =
\begin{bmatrix}
A_{M-1} & A_{M-2} & \dots & A_0 \\
\vdots  & \vdots  & \ddots & \vdots \\
A_{\Ntr-1} & A_{\Ntr-2} & \dots & A_{\Ntr-M}
\end{bmatrix}
\begin{bmatrix} h_0 \\ \vdots \\ h_{M-1} \end{bmatrix}
\quad\Rightarrow\quad \mathbf y = \mathbf A \mathbf h ,
\]
\[
\boxed{\mathbf h_{\mathrm{LS}} = (\mathbf A^H\mathbf A)^{-1}\mathbf A^H \mathbf y}.
\]

\onefig[0.75]{p3_ls_channel}{True discrete equivalent channel vs LS estimate}

Η εκτίμηση ταυτίζεται με το πραγματικό διακριτό ισοδύναμο κανάλι, με σχετικό σφάλμα
$\|\mathbf h_{\mathrm{LS}} - \mathbf h\| / \|\mathbf h\| = 8.3\cdot 10^{-4}$. Αφού το κανάλι είναι
αθόρυβο, το μικρό σφάλμα οφείλεται μόνο στις ουρές του $h(t)$ έξω από τα $M$ δείγματα.

\subsection*{3.6}
\zit{Να υπολογίσετε τον ZF ισοσταθμιστή, μήκους $L \approx 5M$, για καθυστέρηση $\delta = M$.}

Με $L = 50$ και $\delta = 10$, από τις (4.17)–(4.18) ζητάμε $\mathbf H \mathbf f \approx \mathbf e_\delta$,
όπου $\mathbf H$ ο $(M+L-1)\times L$ πίνακας συνέλιξης του $\mathbf h_{\mathrm{LS}}$:
\[
\boxed{\mathbf f_{\mathrm{ZF}} = (\mathbf H^H\mathbf H)^{-1}\mathbf H^H \mathbf e_\delta}.
\]

\onefig[0.7]{p3_zf_total}{Overall response channel - ZF equalizer}

Η συνολική απόκριση προσεγγίζει το $\delta_{n-\delta}$: $|g_\delta| = 0.961$, ενώ η μεγαλύτερη
υπολειπόμενη τιμή είναι $0.148$ και βρίσκεται πριν από τη θέση $\delta$. Ο ZF προσπαθεί να
προσεγγίσει το αντίστροφο του καναλιού με FIR πεπερασμένου μήκους. Επειδή ο ισχυρότερος
συντελεστής είναι ο μεταγενέστερος ($|c_1| > |c_0|$), το αντίστροφο έχει σημαντικό μη αιτιατό
κομμάτι. Η καθυστέρηση $\delta$ «χωράει» μόνο ένα μέρος του, και γι' αυτό η υπολειπόμενη ISI
εμφανίζεται πριν από το $\delta$.

\subsection*{3.7}
\zit{Να φιλτράρετε την symbol-spaced ακολουθία εξόδου, και να αποκόψετε τα πρώτα $\delta$ δείγματα
στην έξοδο του ισοσταθμιστή. Να σχεδιάσετε τα επόμενα $N$ δείγματα. Τι παρατηρείτε;}

Υπολογίζουμε $z_n = \sum_k f_k Y_{n-k}$ και κρατάμε τα $z_\delta, \dots, z_{\delta+N-1}$, τα οποία
αντιστοιχούν στα $A_0,\dots,A_{N-1}$.

\onefig[0.62]{p3_equalized}{Output of the ZF equalizer}

Η ISI έχει αναιρεθεί σε μεγάλο βαθμό: τα δείγματα συγκεντρώνονται γύρω από τα σημεία του
αστερισμού και δεν έχουμε κανένα λάθος (0/200), με $\mathrm{MSE} = 7.2\cdot10^{-2}$. Αφού το
σύστημα είναι αθόρυβο, όλο το «άπλωμα» οφείλεται στην υπολειπόμενη ISI του ZF πεπερασμένου μήκους
που είδαμε στο 3.6.

% =====================================================================
\section{Τέταρτο Μέρος}

\zit{Να επαναλάβετε όλα τα βήματα του 3ου μέρους, μόνο που τώρα να υποθέσετε και CFO.}

Χρησιμοποιούμε το ίδιο πακέτο και το ίδιο φυσικό κανάλι με το τρίτο μέρος, ώστε η σύγκριση να
είναι δίκαιη, και προσθέτουμε CFO $\Delta F_s = 10^{-3}$ ($\Delta f = 0.01$) με τυχαία φάση
$\phi = 4.0186$ rad.

\subsection*{Συγχρονισμός}

\twofig{p4_Cd}{Training statistic (non-ideal channel, CFO)}{p4_Ed}{Channel energy $\mathrm{E}_d$ with CFO}

Το $|\hat h_d|$ εξασθενεί λόγω του $|\alpha|$, όπως στο δεύτερο μέρος, και ο ασθενέστερος λοβός στο
$6T$ σχεδόν χάνεται μέσα στους όρους $n\neq k$. Ο ισχυρός λοβός στο $9.2T$ παραμένει καθαρός.
Σύμφωνα με τις (5.30)–(5.31), επειδή το $|\alpha|$ είναι κοινό για όλα τα $d$, δεν χρειάζεται
να το γνωρίζουμε για τον υπολογισμό του $d^*$.

Βρίσκουμε $d^* = 12$ αντί για $32$ του τρίτου μέρους. Και τα δύο έχουν την ίδια φάση δειγματοληψίας
($d \bmod 10 = 2$) και, όπως είδαμε στο 3.5, το $\mathrm{E}_d$ του πραγματικού καναλιού είναι ίδιο
και στα δύο. Ο συγχρονισμός είναι λοιπόν σωστός. Απλά το παράθυρο των $M$ συντελεστών ξεκινά δύο
σύμβολα νωρίτερα, οπότε οι ισχυροί συντελεστές είναι τώρα οι $h_5$ και $h_8$.

\twofig{p4_output}{Symbol-spaced output (non-ideal channel, CFO)}{p4_equalized_noCFO}{ZF output without CFO correction}

Η ακολουθία εξόδου δεν έχει καμία δομή, αφού στην ISI προστίθεται και η περιστροφή λόγω CFO. Αν
αγνοήσουμε το CFO και εφαρμόσουμε κατευθείαν τα βήματα του τρίτου μέρους (δεξί σχήμα), έχουμε
155 λάθη στα 200 σύμβολα: το μοντέλο $\mathbf y = \mathbf A\mathbf h$ δεν ισχύει όταν η φάση αλλάζει
με τον χρόνο, οπότε και η εκτίμηση του καναλιού και ο ισοσταθμιστής είναι λάθος.

\subsection*{Ταυτόχρονη εκτίμηση καναλιού και \textenglish{CFO}}

Ακολουθούμε το εδάφιο 5.3.1. Η σχέση εισόδου-εξόδου είναι (5.33)
\[
Y_k = e^{j2\pi\Delta f k} \sum_{m=0}^{M-1} h_m A_{k-m},
\]
άρα για τα δείγματα εκπαίδευσης $\mathbf y = \boldsymbol\Gamma(\Delta f)\mathbf A\mathbf h$ και λύνουμε το
πρόβλημα (5.34) $F(\Delta f, \mathbf h) = \|\mathbf y - \boldsymbol\Gamma(\Delta f)\mathbf A\mathbf h\|_2^2$.
Για κάθε $\Delta f$ η βέλτιστη λύση για το $\mathbf h$ είναι η (5.37). Αντικαθιστώντας, με
$\mathbf P = \mathbf A(\mathbf A^H\mathbf A)^{-1}\mathbf A^H$, λαμβάνουμε
\[
\boxed{\Delta f^* = \arg\max_{\Delta f}\; \mathbf y^H \boldsymbol\Gamma(\Delta f)\mathbf P
\boldsymbol\Gamma(\Delta f)^H \mathbf y}, \qquad
\boxed{\mathbf h^* = (\mathbf A^H\mathbf A)^{-1}\mathbf A^H \boldsymbol\Gamma(\Delta f^*)^H \mathbf y}.
\]
Υπολογίζουμε τη συνάρτηση σε πλέγμα $[-\tfrac12, \tfrac12)$ με βήμα $10^{-5}$. Στη διαγώνιο του
$\boldsymbol\Gamma$ χρησιμοποιούμε τους πραγματικούς χρονικούς δείκτες $k = M-1,\dots,\Ntr-1$ των
εξισώσεων (και όχι $0,1,\dots$). Έτσι η φάση του $\mathbf h^*$ συμφωνεί με τη διόρθωση (5.42), η
οποία ξεκινά από το $k = 0$. Διαφορετικά θα έμενε μια σταθερή στροφή
$e^{j2\pi\Delta f(M-1)} \approx 32^\circ$.

\onefig[0.75]{p4_cfo_cost}{Cost function of the joint LS CFO and channel estimation}

Η συνάρτηση έχει ένα καθαρό μέγιστο στο $\Delta f^* = 0.01002$ (πραγματικό $0.01$), δηλαδή
$\Delta\hat F_s = 1.002\cdot 10^{-3}$.

\subsection*{Διόρθωση \textenglish{CFO}, εκτίμηση καναλιού και ισοστάθμιση}

Διορθώνουμε το CFO με (5.42): $\tilde Y_k = Y_k e^{-j2\pi\Delta f^* k}$, $k = 0,\dots,N+M-2$. Κατόπιν
εκτιμούμε το κανάλι LS από την $\tilde Y_k$ (βήμα 3.5), το οποίο δίνει ακριβώς το $\mathbf h^*$
($\|\mathbf h_{\mathrm{LS}} - \mathbf h^*\| = 0$), και υπολογίζουμε τον ZF με $L = 50$, $\delta = 10$.

\twofig{p4_cfo_corrected}{Symbol-spaced output after CFO correction}{p4_ls_channel}{Magnitude of the discrete equivalent channel}

Μετά τη διόρθωση του CFO εμφανίζεται ξανά η δομή των 16 ομάδων του τρίτου μέρους, στραμμένη
κατά μία κοινή φάση. Τα μέτρα των συντελεστών του εκτιμημένου καναλιού ταυτίζονται με αυτά του
πραγματικού διακριτού καναλιού (χωρίς CFO). Η φάση τους διαφέρει κατά την ίδια περίπου γωνία
($\approx -135^\circ$) σε όλους τους σημαντικούς συντελεστές. Η σταθερή φάση $e^{j\theta}$ της
(5.32) «μπαίνει» μέσα στο κανάλι, όπως αναφέρουν οι σημειώσεις.

\twofig{p4_zf_total}{Overall response channel - ZF equalizer (CFO)}{p4_equalized}{ZF output after CFO correction}

Και εδώ δεν έχουμε κανένα λάθος (0/200), αλλά το $\mathrm{MSE} = 0.32$ είναι μεγαλύτερο από το
$0.072$ του τρίτου μέρους και η υπολειπόμενη ISI φαίνεται στον αστερισμό. Η αιτία είναι το
παράθυρο του συγχρονισμού και όχι η εκτίμηση του CFO. Με $d^* = 12$ ο ισχυρός συντελεστής είναι ο
$h_8$ αντί για τον $h_6$, οπότε για σταθερή καθυστέρηση $\delta = M$ μένει λιγότερος χώρος για το
μη αιτιατό κομμάτι του αντιστρόφου (βλ. 3.6). Για επιβεβαίωση, αν η ακολουθία κατασκευαστεί από το
$d = 32$ με το ίδιο CFO, παίρνουμε $\Delta f^* = 0.01000$ και $\mathrm{MSE} = 7.3\cdot10^{-2}$, όσο
περίπου και χωρίς CFO.

% =====================================================================
\section*{Συμπεράσματα}

\begin{itemize}
  \item Ο συγχρονισμός με την ενέργεια βρίσκει σωστά τη φάση δειγματοληψίας και δεν επηρεάζεται
        από το $c$ ή το CFO, αλλά η αρχή του πακέτου δεν ξεχωρίζει καθαρά από τις γειτονικές θέσεις.
  \item Ο συγχρονισμός με σύμβολα εκπαίδευσης δίνει και εκτίμηση του υπερδειγματοληπτημένου
        καναλιού. Με CFO εξασθενεί κατά $|\alpha| = \sigma_A^2|\sin(\pi\Delta f\Ntr)/\sin(\pi\Delta f)|$,
        που επιβεβαιώθηκε αριθμητικά.
  \item Σε ιδανικό κανάλι αρκεί η εκτίμηση CFO με FFT και η εκτίμηση ενός συντελεστή. Σε μη-ιδανικό
        κανάλι χρειάζεται ταυτόχρονη εκτίμηση CFO και καναλιού, αφού χωρίς διόρθωση του CFO ο
        ισοσταθμιστής αποτυγχάνει (155/200 λάθη).
  \item Σε όλα τα τελικά αποτελέσματα ανακτήσαμε το πακέτο χωρίς λάθη. Η ποιότητα του ZF εξαρτάται
        από τη θέση του καναλιού μέσα στο παράθυρο των $M$ συντελεστών σε σχέση με την καθυστέρηση $\delta$.
\end{itemize}

\end{document}
